\section{Proofs}\label{sec:proof}

\subsection{Definitions, assumptions, and auxiliary lemmas}
\begin{figure*}
\begin{minipage}{\textwidth}%
\noindent\fbox{%
    \parbox{\textwidth}
    {{\bf{Notation:}} \\
     $p(x,t)$: distribution on $\cX \times \{0,1\}$\\
     $u = p(t=1)$: the marginal probability of treatment.\\
     $\pt(x) = p(x|t=1)$: treated distribution. $\pc(x) = p(x|t=0)$: control distribution.\\
     $\Phi$: representation function mapping from $\cX$ to  $\cR$.\\
     $\Psi$: the inverse function of $\Phi$, mapping from $\cR$ to $\cX$.\\
     $p_\Phi(r,t)$: the distribution induced by $\Phi$ on $\cR \times \{0,1\}$.\\
       $\pt_\Phi(r)$, $\pc_\Phi(r)$: treated and control distributions induced by $\Phi$ on $\cR$.\\
     $L(\cdot,\cdot)$: loss function, from $\cY \times \cY$ to $\R_+$.\\
     $\lyth$: the expected loss of $h(\Phi(x),t)$ for the unit $x$ and treatment $t$.\\
     $\epsilon_F(h,\Phi)$, $\epsilon_{CF}(h,\Phi)$: expected factual and counterfactual loss of $h(\Phi(x),t)$.\\
     $\tau(x) := \E\left[Y_1 - Y_0 | x \right]$, the expected treatment effect for unit $x$.\\
     $\epehe(f)$: expected error in estimating the individual treatment effect of a function $f(x,t)$. \\
     $\text{IPM}_\cF(p,q)$: the integral probability metric distance induced by function family $\cF$ between distributions $p$ and $q$.
      %$\text{UIPM}_\cF(u_p \cdot p,u_q \cdot q)$: the unnormalized integral probability metric distance induced by function family $\cF$ between distributions $p$ and $q$ scaled respectively by $u_p, u_q \in \R_+$
   \vskip 0pt
    }}
\end{minipage}
\end{figure*}


We first define the necessary distributions and prove some simple results about them. We assume a joint distribution function $p(x,t,Y_0,Y_1)$, such that $(Y_1, Y_0)\indep t | x $, and $0<p(t=1|x)<1$ for all $x$. 
Recall that we assume \emph{Consistency}, that is we assume that we observe $y=Y_1|(t=1)$ and $y=Y_0|(t=0)$.

\begin{thmappdef}\label{def:teA}
The treatment effect for unit $x$ is:
$$\tau(x) := \E\left[Y_1 - Y_0 | x \right].$$
\end{thmappdef}

We first show that under consistency and strong ignorability, the ITE function $\tau(x)$ is identifiable:
\begin{thmapplem}
We have:
\begin{align}
&\E\left[Y_1-Y_0|x\right] =\nonumber \\
&\E\left[Y_1 |x\right]-\E\left[Y_0|x\right] =  \label{eq:useig}\\
&\E\left[Y_1|x,t=1\right]-\E\left[Y_0|x, t=0\right] = \label{eq:useconsistency}\\
&\E\left[y|x,t=1\right]-\E\left[y|x, t=0\right] .\nonumber 
\end{align}
Equality \eqref{eq:useig} is because we assume that $Y_t$ and $t$ are independent conditioned on $x$.
Equality \eqref{eq:useconsistency} follows from the consistency assumption.
Finally, the last equation is composed entirely of observable quantities and can be estimated from data since we assume $0<p(t=1|x)<1$ for all $x$.
\end{thmapplem}
%Recall that we call the marginal distribution over $\cX \times \{0,1\}$ the \emph{factual distribution}, denoted $p^F(x,t)$.
%\begin{thmappdef}\label{def:fcfA}
%Let the counterfactual distribution $p^{CF}(x,t)$ be the distribution over $\cX \times \{0,1\}$ such that $p^{CF}(1-t,x) := p^{F}(t,x)$, where $\cX \subset \R^d$.
%\end{thmappdef}
%An immediate consequence is that
%\begin{equation}\label{covshiftA}
%p^F(x) = p^{CF}(x).
%\end{equation}



\begin{thmappdef}
Let $\pt(x) := p(x|t=1)$, and $\pc(x) := p(x|t=0)$ denote respectively the treatment and control distributions.
\end{thmappdef}



Let $\Phi : \cX \rightarrow \cR$ be a \emph{representation function}. We will assume that $\Phi$ is differentiable.

\begin{thmappasmp}\label{asmp:invA}
The representation function $\Phi $ is one-to-one. Without loss of generality we will assume that $\cR$ is the image of $\cX$ under $\Phi$, and define $\Psi :\cR \rightarrow \cX$ to be the inverse of $\Phi$, such that $\Psi(\Phi(x)) = x$ for all $x \in \cX$.
\end{thmappasmp}

\begin{thmappdef}\label{def:pphiA}
For a representation function $\Phi :\cX \rightarrow \cR$, and for a distribution $p$ defined over $\cX$, let $p_\Phi$ be the distribution induced by $\Phi$ over $\cR$. Define $\pt_\Phi(r) := p_\Phi(r|t=1)$, $\pc_\Phi(r) := p_\Phi(r|t=0)$, to be the treatment and control distributions induced over $\cR$.
\end{thmappdef}


For a one-to-one $\Phi$, the distribution $p_\Phi$ over $\cR \times \{0,1\}$ can be obtained by the standard change of variables formula, using the determinant of the Jacobian of $\Psi(r)$. See \cite{ben1999change} for the case of a mapping $\Phi$ between spaces of different dimensions.

%\begin{thmapplem}\label{lemma:smpl_alg1}
%$p^F_\Phi(r) = p^{CF}_\Phi(r)$
%\end{thmapplem}
%\begin{proof}
%Let $J_\Psi(r)$ be the absolute of the determinant of the Jacobian of $\Psi(r)$. Then by the change of variable formula we have:
%$$p^F_\Phi(r) = J_\Psi(r) \cdot p^F(\Psi(r)) = J_\Psi(r) \cdot p^{CF}(\Psi(r)) = p^{CF}_\Phi(r),$$
%since by \eqref{covshiftA} we have $p^{F}(\Psi(r)) = p^{CF}(\Psi(r))$.
%\end{proof}

\begin{thmapplem}\label{lemma:smpl_alg3}
For all $r \in \cR$, $t \in \{0,1\}$:
\begin{align*}
&p_\Phi(t|r) = p(t|\Psi(r))\\
& p(Y_t|r) = p(Y_t|\Psi(r)).
\end{align*}
\end{thmapplem}
\begin{proof}
Let $J_\Psi(r)$ be the absolute of the determinant of the Jacobian of $\Psi(r)$.
\begin{align*}
&p_\Phi(t|r) = \frac{p_\Phi(t,r)}{p_\Phi(r)}  \stackrel{(a)}{=}
\frac{p(t,\Psi(r))J_\Psi(r)}{p(\Psi(r)) J_\Psi(r)} = \\
&\frac{p(t,\Psi(r))}{p(\Psi(r))} = p(t|\Psi(r)),
\end{align*}
where equality (a) is by the change of variable formula. The proof is identical for $p(Y_t|r)$.
\end{proof}

%\begin{thmapplem}\label{lemma:smpl_alg2}
%For all $r \in \cR$:
%$$p^{CF}_\Phi(r,t) = p^F_\Phi(r|1-t) p^F_\Phi(1-t)$$
%\end{thmapplem}
%\begin{proof}
%$$p^{CF}_\Phi(r,t) = p^{CF}_\Phi(r) p^{CF}_\Phi(t|r) = p^{F}_\Phi(r) p^{F}_\Phi(1-t|r),$$
%since by Lemma \ref{lemma:smpl_alg1} we have $p^{CF}_\Phi(r) = p^{F}_\Phi(r)$, and by Definition \ref{def:fcfA} and Lemma \ref{lemma:smpl_alg3}
%\end{proof}

%%%%%%%%%%%%%%%%%%%%%%%%


Let $L : \cY \times \cY \rightarrow \R_+$ be a loss function, e.g. the absolute loss or squared loss.
\begin{thmappdef}\label{def:perunitlossA}
Let $\Phi : \cX \rightarrow \cR$ be a representation function.
Let $h :\cR \times \{0,1\} \rightarrow \cY$ be an hypothesis defined over the representation space $\cR$.
The expected loss for the unit and treatment pair $(x,t)$ is:
$$\lyth = \int_\cY L(Y_t,h(\Phi(x),t)) p(Y_t|x) dY_t$$
\end{thmappdef}

\begin{thmappdef}\label{def:epfepcf}
The expected factual loss and counterfactual losses of $h$ and $\Phi$ are, respectively:
$$\epsilon_F(h,\Phi) = \int_{\cX \times \{0,1\}} \!\!\!\!\!\!\! \lyth\, p(x,t)\, dxdt $$
$$\epsilon_{CF}(h,\Phi) = \int_{\cX \times \{0,1\}} \!\!\!\!\!\!\! \lyth\, p(x,1-t)\, dxdt .$$
\end{thmappdef}
When it is clear from the context, we will sometimes use $\epsilon_F(f)$ and $\epsilon_{CF}(f)$ for the expected factual and counterfactual losses of an arbitrary function $f: \cX \times \{0,1\} \rightarrow \cY$.

\begin{thmappdef}\label{def:decomploss}
The expected \emph{treated} and \emph{control} losses are:
$$\epsilon^{t=1}_F(h,\Phi) = \int_{\cX } \!\!\! \lyxoneh\, \pt(x)\, dx $$
$$\epsilon^{t=0}_F(h,\Phi) = \int_{\cX } \!\!\! \lyxzeroh\, \pc(x)\, dx $$
$$\epsilon^{t=1}_{CF}(h,\Phi) = \int_{\cX} \!\!\! \lyxoneh\, \pc(x)\, dx $$
$$\epsilon^{t=0}_{CF}(h,\Phi) = \int_{\cX } \!\!\! \lyxzeroh\, \pt(x)\, dx .$$
\end{thmappdef}
The four losses above are simply the loss conditioned on either the control or treated set.  Let $u := p(t=1)$ be the proportion of treated in the population. We then have the immediate result:
\begin{thmapplem}\label{lem:epsdecopm}
$$\epsilon_F(h,\Phi) = u\cdot \epsilon^{t=1}_F(h,\Phi) + (1-u)\cdot \epsilon^{t=0}_F(h,\Phi)$$
$$\epsilon_{CF}(h,\Phi) = (1-u)\cdot \epsilon^{t=1}_{CF}(h,\Phi) + u\cdot  \epsilon^{t=0}_{CF}(h,\Phi).$$
\end{thmapplem}
The proof is immediate, noting that $p(x,t) = u\cdot \pt(x) + (1-u)\cdot \c(x)$, and from the Definitions \ref{def:perunitlossA} and \ref{def:decomploss} of the losses.

\begin{thmappdef}\label{def:ipm}
Let $\cF$ be a function family consisting of functions $g: \cS\rightarrow \R$. For a pair of distributions $p_1$, $p_2$ over $\cS$, define the \emph{Integral Probability Metric}:
$$\text{IPM}_\cF(p_1,p_2) = \sup_{g\in \cF} \left|\int_\cS g(s) \left(p_1(s) - p_2(s)\right) \, ds \right|$$
\end{thmappdef}
$\text{IPM}_\cF(\cdot,\cdot)$ defines a pseudo-metric on the space of probability functions over $\cS$, and for sufficiently large function families, $\text{IPM}_\cF(\cdot,\cdot)$ is a proper metric \cite{muller1997integral}. Examples of sufficiently large functions families includes the set of bounded continuous functions, the set of $1$-Lipschitz functions, and the set of unit norm functions in a universal Reproducing Norm Hilbert Space. The latter two give rise to the Wasserstein and Maximum Mean Discrepancy metrics, respectively \cite{gretton2012mmd,sriperumbudur2012empirical}. We note that for function families $\cF$ such as the three mentioned above, for which $g \in \cF \implies -g \in \cF$, the absolute value can be omitted from definition \ref{def:ipm}.


\subsection{General IPM bound}


We now state and prove the most important technical lemma of this section.
\begin{thmapplem}[Lemma 1, main text]\label{lem:genA}
Let $\Phi: \cX \rightarrow \cR$ be an invertible representation with $\Psi$ its inverse. Let $\pt_\Phi, \pc_\Phi$ be defined as in Definition \ref{def:pphiA}. Let $u = p(t=1)$. Let $\cF$ be a family of functions $g:\cR \rightarrow \R$, and denote by $\text{IPM}_\cF(\cdot, \cdot)$ the integral probability metric induced by $\cF$.  Let $h : \cR \times \{0,1\} \rightarrow \cY$ be an hypothesis. Assume there exists a constant $B_\Phi>0$, such that for $t=0,1$, the function $g_{\Phi,h}(r,t) := \frac{1}{B_\Phi} \cdot  \lythr \in \cF$. Then we have:
\begin{align}
&\epsilon_{CF}(h,\Phi) \leq \nonumber \\
&(1-u)  \epsilon^{t=1}_F(h,\Phi) + u  \epsilon^{t=0}_F(h,\Phi) + \nonumber \\ 
&  B_\Phi \cdot \text{IPM}_\cF\left( \pt_\Phi, \pc_\Phi \right).
\end{align}
\end{thmapplem}

\begin{proof}
\allowdisplaybreaks
\begin{align}
&\epsilon_{CF}(h,\Phi) - \left[ (1-u)\cdot \epsilon^{t=1}_{F}(h,\Phi) + u\cdot  \epsilon^{t=0}_{F}(h,\Phi) \right]= \nonumber \\
&\left[ (1-u)\cdot \epsilon^{t=1}_{CF}(h,\Phi) + u\cdot  \epsilon^{t=0}_{CF}(h,\Phi)\right] - \nonumber \\ 
&\left[ (1-u)\cdot \epsilon^{t=1}_{F}(h,\Phi) + u\cdot  \epsilon^{t=0}_{F}(h,\Phi)\right] = \nonumber  \\
&(1-u)\cdot \left[\epsilon^{t=1}_{CF}(h,\Phi)  - \epsilon^{t=1}_{F}(h,\Phi)\right]  + \nonumber \\
&u\cdot \left[\epsilon^{t=0}_{CF}(h,\Phi)  - \epsilon^{t=0}_{F}(h,\Phi)\right]  = \label{eq:usedef6}\\
&(1-u)\int_{\cX} \!\!\! \lyxoneh\, \left(\pc(x) - \pt(x)\right) \, dx  + \nonumber \\
& u  \int_{\cX} \!\!\!  \lyxzeroh\, \left(\pt(x) - \pc(x) \right)\, dx = \label{eq:usecov} \\
&(1-u)\int_{\cR} \!\!\! \lyonehpsi\, \left(\pc_\Phi(r) - \pt_\Phi(r)\right) \, dr  + \nonumber \\
& u  \int_{\cR} \!\!\!  \lyzerohpsi\, \left(\pt_\Phi(r) - \pc_\Phi(r) \right)\, dr = \nonumber \\
&B_\Phi \cdot (1-u)\int_{\cR} \!\! \frac{1}{B_\Phi}\lyonehpsi\, \left(\pc_\Phi(r) - \pt_\Phi(r)\right) \, dr  + \nonumber \\
& B_\Phi \cdot  u  \int_{\cR} \!\! \frac{1}{B_\Phi} \lyzerohpsi\, \left(\pt_\Phi(r) - \pc_\Phi(r) \right)\, dr \leq \label{proof:ineq} \\
&B_\Phi \cdot (1-u)\sup_{g\in \cF} \left| \int_{\cR} \!\! g(r)\, \left(\pc_\Phi(r) - \pt_\Phi(r)\right) \, dr  \right| + \nonumber \\
& B_\Phi \cdot u \sup_{g\in \cF} \left| \int_{\cR} \!\! g(r)\, \left(\pt_\Phi(r) - \pc_\Phi(r) \right)\, dr \right| = \label{eq:useipmdef}\\
&B_\Phi \cdot \text{IPM}_\cF(\pc_\Phi,\pt_\Phi) .
\end{align}
Equality \eqref{eq:usedef6} is by Definition \ref{def:decomploss} of the treated and control loss, equality \eqref{eq:usecov} is by the change of variables formula and Definition \ref{def:pphiA} of $\pt_\Phi$ and $\pc_\Phi$, inequality \eqref{proof:ineq} is by the premise that $\frac{1}{B_\Phi} \cdot  \lythr \in \cF$ for $t=0,1$, and \eqref{eq:useipmdef} is by Definition \ref{def:ipm} of an IPM.
\end{proof}


The essential point in the proof of Lemma \ref{lem:genA} is inequality \ref{proof:ineq}. Note that on the l.h.s. of the inequality, we need to evaluate the expectations of $ \lyzerohpsi$ over $\pt_\Phi$ and $\lyonehpsi$ over $\pc_\Phi$. Both of these expectations are in general unavailable, since they require us to evaluate treatment outcomes on the control, and control outcomes on the treated. We therefore upper bound these unknowable quantities by taking a supremum over a function family which includes $\lyzerohpsi$ and $\lyonehpsi$. The upper bound ignores most of the details of the outcome, and amounts to measuring a distance between two distributions we have samples from: the control and treated distribution. Note that for a randomized trial (i.e. when $t \indep x$) with we have that $\text{IPM}(\pt_\Phi,\pc_\Phi) = 0$. Indeed, it is straightforward to show that in that case we actually have an equality: $\epsilon_{CF}(h,\Phi) = (1-u)\cdot \epsilon^{t=1}_{F}(h,\Phi) + u\cdot  \epsilon^{t=0}_{F}(h,\Phi)$.


The crucial condition in Lemma \ref{lem:genA} is that the function $g_{\Phi,h}(r) := \frac{1}{B_\Phi} \lythr$ is in $\cF$. In subsections \ref{subsec:wass} and \ref{subsec:mmd} below we look into two specific function families $\cF$, and evaluate what does this inclusion condition entail, and in particular we will derive specific bounds for $B_\Phi$.



\begin{thmappdef}\label{def:m}
For $t=0,1$ define:
$$m_t(x) := \E\left[Y_t|x\right].$$
\end{thmappdef}
Obviously for the treatment effect $\tau(x)$ we have $\tau(x) = m_1(x)-m_0(x)$.


Let $f: \cX \times \{0,1\} \rightarrow \cY$ by an hypothesis, such that $f(x,t) = h(\Phi(x),t)$ for a representation $\Phi$ and hypothesis $h$ defined over the output of $\Phi$.
\begin{thmappdef}\label{def:inditeerrA}
The treatment effect estimate for unit $x$ is:
\begin{align*}
\hat{\tau}_f  (x) = f(x,1) - f(x,0).
\end{align*}
\end{thmappdef}


\begin{thmappdef}
 The expected Precision in Estimation of Heterogeneous Effect (PEHE) loss of $g$ is:
$$\epehe(f) = \int_{\cX } \left( \hat{\tau}_f(x) - \tau(x) \right)^2 \, p(x) \, dx .$$
\end{thmappdef}


\begin{thmappdef}\label{def:sigma2}
The expected variance of $Y_t$ with respect to a distribution $p(x,t)$:
$$\sigma^2_{Y_t}(p(x,t)) = \int_{\cX \times \cY} \left(Y_t - m_t(x)\right)^2 p(Y_t|x)p(x,t) \, dY_t dx .$$
We define:
\begin{align*}
&\sigma^2_{Y_t} = \min\{\sigma^2_{Y_t}(p(x,t)), \sigma^2_{Y_t}(p(x,1-t))\} , \\
&\sigma^2_{Y} = \min\{\sigma^2_{Y_0},\sigma^2_{Y_1}\}.
\end{align*}
\end{thmappdef}
If $Y_t$ are deterministic functions of $x$, then $\sigma^2_{Y} = 0$.


We now show that $\epehe(f)$ is upper bounded by $2 \epsilon_F + 2 \epsilon_{CF} - 2\sigma^2_Y$ where $\epsilon_F$ and $\epsilon_{CF}$ are w.r.t. to the squared loss. An analogous result can be obtained for the absolute loss, using mean absolute deviation. 

\begin{thmapplem}\label{lem:epsmloss}
For any function $f : \cX \times \{0,1\} \rightarrow \cY$, and distribution $p(x,t)$ over $\cX \times \{0,1\}$:
\begin{align*}
\int_\cX& \left(f(x,t) - m_t(x) \right)^2 \, p(x,t) \, dx dt = \\
&\epsilon_F(f) - \sigma^2_{Y_t}(p(x,t)), \\
 \int_\cX &\left(f(x,t) - m_t(x) \right)^2 \, p(x,1-t) \, dx dt = \\
& \epsilon_{CF}(f) - \sigma^2_{Y_t}(p(x,1-t)),
\end{align*}
where $\epsilon_F(f)$ and $\epsilon_{CF}(f)$ are w.r.t. to the squared loss.
\end{thmapplem}
\begin{proof}
For simplicity we will prove for $p(x,t)$ and $\epsilon_F(f)$. The proof for $p(x,1-t)$ and $\epsilon_{CF}$ is identical.
\begin{align}
&\epsilon_F(f) = \nonumber \\
&\int_{\cX  \times \{0,1\} \times \cY}   \!\!\!\!\!\!\!\!\!\!\!\!\!\!\!\!\!\!\!\!\!\! \left(f(x,t) - Y_t\right)^2 p(Y_t|x) p(x,t) \, dY_t dx dt = \nonumber \\
&\int_{\cX  \times \{0,1\} \times \cY}   \!\!\!\!\!\!\!\!\!\!\!\!\!\!\!\!\!\!\!\!\!\! \left(f(x,t) - m_t(x)\right)^2 p(Y_t|x) p(x,t) \, dY_t dx dt + \nonumber \\
&\int_{\cX  \times \{0,1\} \times \cY}   \!\!\!\!\!\!\!\!\!\!\!\!\!\!\!\!\!\!\!\!\!\!  \left(m_t(x) - Y_t\right)^2 p(Y_t|x) p(x,t) \, dY_t dx dt + \label{eq:expctzero} \\
&\int_{\cX  \times \{0,1\} \times \cY}  \!\!\!\!\!\!\!\!\!\!\!\!\!\!\!\!\!\!\!\!\!\! \left(f(x,t) - m_t(x)\right)\left(m_t(x) - Y_t\right) p(Y_t|x) p(x,t) \, dY_t dx dt = \label{eq:epssigm} \\
&\int_{\cX  \times \{0,1\} }   \!\!\!\!\!\!\!\!\!\!\!\!\!\!\!\!\!\! \left(f(x,t) - m_t(x)\right)^2  p(x,t) \,  dx dt + \nonumber \\
&\quad \quad \sigma^2_{Y_0}(p(x,t)) + \sigma^2_{Y_1}(p(x,t)) +0 , \nonumber
\end{align}
where the equality \eqref{eq:epssigm} is by the Definition \ref{def:sigma2} of $\sigma^2_{Y_t}(p)$, and because the integral in \eqref{eq:expctzero} evaluates to zero, since $m_t(x) = \int_\cX Y_t p(Y_t|x) \, dx$.
\end{proof}

\begin{thmappthm}\label{thm:indtausqloss}
Let $\Phi : \cX \rightarrow \cR$ be a one-to-one representation function, with inverse $\Psi$. Let $\pt_\Phi, \pc_\Phi$ be defined as in Definition \ref{def:pphiA}. Let $u = p(t=1)$. Let $\cF$ be a family of functions $g:\cR \rightarrow \R$, and denote by $\text{IPM}_\cF(\cdot, \cdot)$ the integral probability metric induced by $\cF$.  Let $h : \cR \times \{0,1\} \rightarrow \cY$ be an hypothesis. Let the loss $L(y_1,y_2) = (y_1 - y_2)^2$.  Assume there exists a constant $B_\Phi>0$, such that for $t \in \{0,1\}$, the functions $g_{\Phi,h}(r,t) := \frac{1}{B_\Phi} \cdot  \lythr \in \cF$.
We then have:
\begin{align*}
&\epehe(h,\Phi) \leq  \nonumber \\
& 2\!\left(\epsilon_{CF}(h,\Phi) + \epsilon_F(h,\Phi) - 2\sigma^2_Y \right) \leq  \nonumber \\
&2\!\left(\epsilon_F^{t=0}(h,\Phi)\! +\!\epsilon_F^{t=1}(h,\Phi)\!+ \! B_\Phi  \text{IPM}_\cF \left( \pt_\Phi, \pc_\Phi \right)\! -\! 2\sigma^2_Y\right)\!,
\end{align*}
where $\epsilon_F$ and $\epsilon_{CF}$ are with respect to the squared loss.
\end{thmappthm}

\begin{proof}
We will prove the first inequality, $\epehe(f) \leq   2\epsilon_{CF}(h,\Phi) + 2\epsilon_F(h,\Phi) - 2\sigma^2_Y$. The second inequality is then immediate by Lemma \ref{lem:genA}.
Recall that we denote $\epehe(f) = \epehe(h,\Phi)$ for $f(x,t) = h(\Phi(x),t)$.
%We will bound $\int_{\cX } \left| \hat{\tau}_f(x) - \tau(x) \right| \, p(x) \, dx$.
\begin{align}
&\epehe(f) = \nonumber\\
& \int_\cX \bigl(\left(f(x,1) - f(x,0)\right) - \left(m_1(x) - m_0(x)\right) \bigr)^2 p(x) \, dx= \nonumber\\
& \int_\cX \bigl( \left(f(x,1) - m_1(x)\right) + \left(m_0(x) - f(x,0)\right) \bigr)^2 p(x) \, dx \leq  \label{eq:sqineq}\\
&  2\int_\cX \left(\left(  f(x,1) - m_1(x) \right)^2  + \left(m_0(x) - f(x,0) \right)^2\right) p(x) \, dx = \label{eq:use:margt2sq} \\
& 2\int_\cX  \left(f(x,1) - m_1(x)\right)^2  p(x,t=1)  \,dx+ \nonumber \\
& \quad 2\int_\cX \left(m_0(x) - f(x,0) \right)^2  p(x,t=0)  \, dx +  \nonumber\\
& \quad 2\int_\cX  \left(f(x,1) - m_1(x)\right)^2 p(x,t=0)  \, dx+ \nonumber\\
& \quad 2\int_\cX \left(m_0(x) - f(x,0) \right)^2  p(x,t=1)  \, dx \nonumber = \\
&2 \int_\cX \left(f(x,t) - m_t(x)\right)^2 p(x,t) \, dx dt + \nonumber \\
&\quad 2 \int_\cX \left(f(x,t) - m_t(x)\right)^2 p(x,1-t)\, dx dt \leq \label{eq:usesigmineq}\\
&2 (\epsilon_F - \sigma^2_Y) + 2 (\epsilon_{CF} - \sigma^2_Y). \nonumber
\end{align}
where \eqref{eq:sqineq} is because $(x+y)^2 \leq 2(x^2 + y^2)$, \eqref{eq:use:margt2sq} is because $p(x) = p(x,t=0) + p(x,t=1)$ and \eqref{eq:usesigmineq} is by Lemma \ref{lem:epsmloss} and Definition \ref{def:epfepcf} of the losses $\epsilon_F$, $\epsilon_{CF}$ and Definition \ref{def:sigma2} of $\sigma^2_Y$.
Having established the first inequality in the Theorem statement, we now show the second.
We have by Lemma \ref{lem:genA} that:
\begin{align*}
&\epsilon_{CF}(h,\Phi) \leq  \\ 
&(1-u)  \epsilon^{t=1}_F(h,\Phi) + u  \epsilon^{t=0}_F(h,\Phi) +   B_\Phi \cdot \text{IPM}_\cF\left( \pt_\Phi, \pc_\Phi \right).
\end{align*}
We further have by Lemma \ref{lem:epsdecopm} that:
\begin{align*} 
\epsilon_F(h,\Phi) = u \epsilon^{t=1}_F(h,\Phi) + (1-u) \epsilon^{t=0}_F(h,\Phi).
\end{align*}
Therefore 
\begin{align*}
&\epsilon_{CF}(h,\Phi)+\epsilon_F(h,\Phi) \leq  \\
&\epsilon^{t=1}_F(h,\Phi) + \epsilon^{t=0}_F(h,\Phi)+ B_\Phi  \text{IPM}_\cF\left( \pt_\Phi, \pc_\Phi \right).
\end{align*}
\end{proof}

The upper bound is in terms of the standard generalization error on the treated and control distributions separately. Note that in some cases we might have very different sample sizes for treated and control, and that will show up in the finite sample bounds of these generalization errors.

We also note that the upper bound can be easily adapted to the case of the absolute loss PEHE $\left| \hat{\tau}(x) - \tau(x)\right|$. In that case the upper bound in the Theorem will have a factor $1$ instead of the $2$ stated above, and the standard deviation $\sigma^2_Y$ replaced by mean absolute deviation. The proof is straightforward where one simply applies the triangle inequality in inequality \eqref{eq:sqineq}.


We will now give specific upper bounds for the constant $B_\Phi$ in Theorem \ref{thm:indtausqloss}, using two function families $\cF$ in the IPM: the family of $1$-Lipschitz functions, and the family of $1$-norm reproducing kernel Hilbert space functions. Each one will have different assumptions about the distribution $p(x,t,Y_0,Y_1)$ and about the representation $\Phi$ and hypothesis $h$.

\subsection{The family of $1$-Lipschitz functions}\label{subsec:wass}

For $\cS \subset \R^d$, a function $f : \cS \rightarrow \R$ has Lipschitz constant $K$ if for all $x,y \in \cS$, $|f(x) - f(y)| \leq K\|x-y\|$. If $f$ is differentiable, then a sufficient condition for $K$-Lipschitz constant is if $\| \frac{\partial f}{\partial s} \| \leq K$ for all $s\in \cS$.

For simplicity's sake we assume throughout this subsection that the true labeling functions the densities $p(Y_t|x)$ and the loss $L$ are differentiable. However, this assumption could be relaxed to a mere Lipschitzness assumption.
\begin{thmappasmp}\label{assmp_f_lip}
There exists a constant $K >0$ such that for all $x \in \cX$, $t \in \{0,1\}$, $\| \frac{p(Y_t|x)}{\partial x} \| \leq K$.
\end{thmappasmp}
Assumption \ref{assmp_f_lip} entails that each of the potential outcomes change smoothly as a function of the covariates (context) $x$.

\begin{thmappasmp}\label{asmp:loss}
The loss function $L$ is differentiable, and there exists a constant $K_L >0$ such that $\left|\frac{d L(y_1,y_2)}{d y_i}\right| \leq K_L$ for $i=1,2$. Additionally, there exists a constant $M$ such that for all $y_2\in \cY$, $M \geq \int_{\cY} L(y_1,y_2)\, dy_1$.
\end{thmappasmp}
Assuming $\cY$ is compact, loss functions which obey Assumption \ref{asmp:loss} include the log-loss, hinge-loss, absolute loss, and the squared loss.

When we let $\cF$ in Definition \ref{def:ipm} be the family of $1$-Lipschitz functions, we obtain the so-called \emph{$1$-Wasserstein} distance between distributions, which we denote $\text{Wass}(\cdot,\cdot)$.
It is well known that $\text{Wass}(\cdot,\cdot)$ is indeed a metric between distributions \cite{villani2008optimal}.




\begin{thmappdef}\label{def:rcondA}
Let $\Jphix$ be the Jacobian matrix of $\Phi$ at point $x$, i.e. the matrix of the partial derivatives of $\Phi$. Let $\sigma_{max}(A)$ and $\sigma_{min}(A)$ denote respectively the largest and smallest singular values of a matrix $A$.
Define $\rho(\Phi) = \sup_{x \in \cX} \, \, \sigma_{max}\left(\Jphix\right) / \sigma_{min}\left(\Jphix\right)$.
\end{thmappdef}
It is an immediate result that $\rho(\Phi) \geq 1$.

\begin{thmappdef}
We will call a representation function $\Phi : \cX \rightarrow \cR$ \emph{Jacobian-normalized} if $\sup_{x \in \cX} \sigma_{max} \left(\Jphix \right) = 1$.
\end{thmappdef}
Note that any non-constant representation function $\Phi$ can be Jacobian-normalized by a simple scalar multiplication.


\begin{thmapplem}\label{lem:lip}
Assume that $\Phi$ is a Jacobian-normalized representation, and let $\Psi$ be its inverse.
For $t=0,1$, the Lipschitz constant of $p(Y_t|\Psi(r))$ is bounded by $\rho(\Phi) K$, where  $K$ is from Assumption \ref{assmp_f_lip}, and $\rho(\Phi)$ as in Definition \ref{def:rcondA}.
\end{thmapplem}
\begin{proof}
Let $\Psi: \cR \rightarrow \cX$ be the inverse of $\Phi$, which exists by the assumption that $\Phi$ is one-to-one. Let $\dPhi$ be the Jacobian matrix of $\Phi$ evaluated at $x$, and similarly let $\dPsi$ be the Jacobian matrix of $\Psi$ evaluated at $r$. Note that $\dPsi \cdot \dPhi = I$ for $r= \Phi(x)$, since $\Psi \circ \Phi$ is the identity function on $\cX$.	Therefore for any $r \in \cR$ and $x= \Psi(r)$:
\begin{equation}\label{eqn:sigman2}
\sigma_{max}\left(\dPsi\right) = \frac{1}{\sigma_{min}\left(\dPhi\right)},
\end{equation}
where $\sigma_{max}(A)$ and $\sigma_{min}(A)$ are respectively the largest and smallest singular values of the matrix $A$, i.e. $\sigma_{max}(A)$ is the spectral norm of $A$.


For $x = \Psi(r)$ and $t \in \{0,1\}$, we have by the chain rule:
\begin{align}
& \| \frac{\partial p(Y_t|\Psi(r))}{\partial r} \| =  \| \frac{\partial p(Y_t|\Psi(r))}{\partial \Psi(r)}  \dPsi\|  \leq \label{eqn:matnorm}\\
& \| \dPsi\|  \| \frac{\partial p(Y_t|\Psi(r))}{\partial \Psi(r)} \|  = \label{eqn:invjac} \\
& \frac{1}{\sigma_{min}\left(\dPhi\right)} \| \frac{\partial p(Y_t|x)}{\partial x} \| \leq \label{eqn:maxgrad}\\
&  \frac{K}{\sigma_{min}\left(\dPhi\right)} \leq  \rho(\Phi) K \label{eqn:ineq_vphi},
\end{align}
where inequality \eqref{eqn:matnorm} is by the matrix norm inequality, equality \eqref{eqn:invjac} is by \eqref{eqn:sigman2}, inequality \eqref{eqn:maxgrad} is by assumption \ref{assmp_f_lip} on the norms of the gradient of $p(Y_t|x)$ w.r.t $x$ , and inequality \eqref{eqn:ineq_vphi} is by Definition \ref{def:rcondA} of $\rho(\Phi)$, the assumption that $\Phi$ is Jacobian-normalized, and noting that singular values are necessarily non-negative.

%Inequalities \eqref{eqn:matnorm} - \eqref{eqn:ineq_vphi}, along with Assumption \ref{assmp1} and \ref{lemma:lip_sum}, show us that the Lipschitz constant of %$\tilde{f_t}$ is bounded by $\sqrt{\frac{K^2}{m(\Phi)^2} + b^2 K^2} = K \sqrt{\frac{1}{m(\Phi)^2} + b^2}$.

\end{proof}

\begin{thmapplem}\label{lem:errlip}
Under the conditions of Lemma \ref{lem:genA}, further assume that for $t=0,1$, $p(Y_t|x)$ has gradients bounded by $K$ as in \ref{assmp_f_lip}, that $h$ has bounded gradient norm $b K$, that the loss $L$ has bounded gradient norm $K_L$, and that $\Phi$ is Jacobian-normalized. Then the Lipschitz constant of $ \lythr$ is upper bounded by $K_L \cdot K \left(M \rho(\Phi) +  b \right)$ for $t=0,1$.
\end{thmapplem}
\begin{proof}
Using the chain rule, we have that:
\begin{align}
& \|\frac{\partial \lythr}{\partial r}\| =\| \frac{\partial}{\partial r} \int_\cY L(Y_t,h(r,t))  p(Y_t|r) dY_t \|= \nonumber \\
& \| \int_\cY \frac{\partial}{\partial r} \left[ L(Y_t,h(r,t))  p(Y_t|r)\right] \, dY_t \| = \nonumber \\
& \| \int_\cY \! p(Y_t|r)\frac{\partial}{\partial r} \! L(Y_t,h(r,t)) \!   +\! L(Y_t,h(r,t)) \frac{\partial}{\partial r} \! p(Y_t|r)  dY_t \|  \leq  \nonumber \\
& \int_\cY p(Y_t|r) \| \frac{\partial}{\partial r} L(Y_t,h(r,t)) \| \, dY_t + \nonumber\\
& \int_\cY  L(Y_t,h(r,t)) \frac{\partial}{\partial r} p(Y_t|r) \, dY_t \leq \label{lip_const_ineq1}\\
& \int_\cY p(Y_t|r) \| \frac{\partial L(Y_t,h(r,t))}{\partial h(r,t) }  \frac{\partial h(r,t)}{\partial r} \| \, dY_t + \nonumber\\
& \int_\cY  L(Y_t,h(r,t)) \frac{\partial}{\partial r} p(Y_t|r) \, dY_t \leq \label{lip_const_ineq2}\\
& \int_\cY p(Y_t|r) K_L \cdot b \cdot K + M \cdot \rho(\Phi) \cdot  K,
\end{align}
%\begin{align*}
%&\frac{\partial L}{\partial r} = \frac{\partial L}{\partial Y_t(\Psi(r))} \frac{\partial Y_t(\Psi(r))}{\partial r} + \frac{\partial L}{\partial h(r,t)} \frac{\partial h(r,t)}{\partial r} \leq \\
%&  K_l K \rho(\Phi) + K_l \cdot  b K = K_L \cdot K \left(\rho(\Phi) +  b \right),
%\end{align*}
where inequality \ref{lip_const_ineq1} is due to Assumption \ref{asmp:loss} and inequality \ref{lip_const_ineq2} is due to Lemma \ref{lem:lip}.
\end{proof}



\begin{thmapplem}\label{thm:wasA}
Let $u = p(t=1)$ be the marginal probability of treatment, and assume $0<u<1$. Let $\Phi : \cX \rightarrow \cR$ be a one-to-one, Jacobian-normalized representation function. Let $K$ be the Lipschitz constant of the functions $p(Y_t|x)$ on $\cX$. Let $K_L$ be the Lipschitz constant of the loss function $L$, and $M$ be as in Assumption \ref{asmp:loss}. Let $h : \cR \times \{0,1\} \rightarrow \R$ be an hypothesis with Lipschitz constant $b K$.
%Denote by $\text{Wass}(u \cdot \pt_\Phi,(1-u) \cdot \pc_\Phi)$ the Wasserstein distance between the treatment and control distributions induced by $\Phi$.
Then:
\begin{align}\label{eq:thm_wasA}
&\epsilon_{CF}(h,\Phi) \leq  \nonumber \\
&(1-u)  \epsilon^{t=1}_F(h,\Phi) + u  \epsilon^{t=0}_F(h,\Phi)  + \nonumber \\
& 2 \left(M \rho(\Phi) +b\right) \cdot K  \cdot K_L \cdot \text{Wass}(\pt_\Phi \, , \pc_\Phi).
\end{align}
\end{thmapplem}

\begin{proof}
We will apply Lemma \ref{lem:genA} with $\cF = \{g: \cR \rightarrow \R \text{ s.t. } f  \text{ is } 1 \text{-Lipschitz}\}$. By Lemma \ref{lem:errlip}, we have that for $B_\Phi = \left(M \rho(\Phi) +b\right) \cdot K  \cdot K_L$, the function $\frac{1}{B_\Phi} \lythr \in \cF$. Inequality \eqref{eq:thm_wasA} then holds as a special case of Lemma \ref{lem:genA}.
\end{proof}
\begin{thmappthm}%[Theorem 2, main paper]
Under the assumptions of Lemma \ref{thm:wasA}, using the squared loss for $\epsilon_F$, we have:
\begin{align*}
&\epehe(h,\Phi) \leq \nonumber \\
&2 \epsilon^{t=0}_F(h,\Phi)  + 2\epsilon^{t=1}_F(h,\Phi)  - 4\sigma^2_Y +  \nonumber \\
&2 \left(M \rho(\Phi) +b\right) \cdot K  \cdot K_L \cdot \text{Wass}( \pt_\Phi \, ,  \pc_\Phi) .
\end{align*}
\end{thmappthm}
\begin{proof}
Plug in the upper bound of Lemma \ref{thm:wasA} into the upper bound of Theorem \ref{thm:indtausqloss}.
\end{proof}


We examine the constant $\left(M \rho(\Phi) +b\right) \cdot K \cdot K_L$ in Theorem \ref{thm:wasA}. $K$, the Lipschitz constant of $m_0$ and $m_1$, is not under our control and measures an aspect of the complexity of the true underlying functions we wish to approximate. The terms $K_L$ and $M$ depend on our choice of loss function and the size of the space $\cY$. The term $b$ comes from our assumption that the hypothesis $h$ has norm $ b K$. Note that smaller $b$, while reducing the bound, might force the factual loss term $\epsilon_F(h,\Phi)$ to be larger since a small $b$ implies a less flexible $h$. Finally, consider the term $\rho(\Phi)$.
The assumption that $\Phi$ is normalized is rather natural, as we do not expect a certain scale from a representation. Furthermore, below we show that in fact the Wasserstein distance is positively homogeneous with respect to the representation $\Phi$. Therefore, in Lemma \ref{thm:wasA}, we can indeed assume that $\Phi$ is normalized. The specific choice of \emph{Jacobian-normalized} scaling yields what is in our opinion a more interpretable result in terms of the inverse condition number $\rho(\Phi)$. For twice-differentiable $\Phi$, $\rho(\Phi)$ is minimized if and only if $\Phi$ is a linear orthogonal transformation \cite{mathover}. %\todo[inline]{Fill in details for upper bounding $\rho(\Phi)$.}

\begin{thmapplem}\label{lem:wassinvar}
The Wasserstein distance is positive homogeneous for scalar transformations of the underlying space. Let $p$, $q$ be probability density functions defined over $\cX$. For $\alpha>0$ and
 the mapping $\Phi(x) = \alpha x$, let $p_\alpha$ and $q_\alpha$ be the distributions on $\alpha\cX$ induced by $\Phi$. Then:
$$\text{Wass}\left(p_{\alpha  },q_{\alpha  }\right) = \alpha\text{Wass}\left(p,q\right).$$
\end{thmapplem}

\begin{proof}
Following \cite{villani2008optimal,tabak2016preconditioning}, we use another characterization of the Wasserstein  distance. Let $\mathcal{M}_{p,q}$ be the set of mass preserving maps from $\cX$ to itself which map the distribution $p$ to the distribution $q$. That is,  $\mathcal{M}_{p,q} = \{M : \cX \rightarrow \cX \text{ s.t. }q(M(S)) = p(S) \text{ for all measurable bounded } S \subset \cX\}$. We then have that:
\begin{equation}\label{eq:wass_def_M}
\text{Wass}(p,q) = \inf_{M \in \mathcal{M}_{p,q}} \int_\cX \|M(x) - x \| p(x) \, dx .
\end{equation}
It is known that the infimum in \eqref{eq:wass_def_M} is actually achievable \citep[Theorem 5.2]{villani2008optimal}. Denote by $M^*: \cX \rightarrow \cX$ the map achieving the infimum for $\text{Wass}(p,q)$ . Define $M^*_\alpha :\alpha \cX \rightarrow \alpha \cX$, by $M^*_\alpha(x') = \alpha M^*(\frac{x'}{\alpha})$, where $x' = \alpha x$.  $M^*_\alpha$ maps $p_\alpha$ to $q_\alpha$, and we have that $\|M^*_\alpha(x') - x'\| = \alpha \|M^*(x)-x\|$. Therefore $M^*_\alpha$ achieves the infimum for the pair $(p_\alpha,q_\alpha)$, and we have that $\text{Wass}\left(p_{\alpha  },q_{\alpha  }\right) = \alpha\text{Wass}\left(p,q\right)$.
\end{proof}

\subsection{Functions in the unit ball of a RKHS}\label{subsec:mmd}

Let $\cH_x, \cH_r$ be a reproducing kernel Hilbert space, with corresponding kernels $k_x(\cdot,\cdot)$, $k_r(\cdot,\cdot)$. We have for all $x \in \cX$ that  $k_x(\cdot,x)$ is its Hilbert space mapping, and similarly $k_r(\cdot,r)$ for all $r \in \cR$.

Recall that the major condition in Lemma \ref{lem:genA} is that $\frac{1}{B_\Phi} \lythr \in \cF$. The function space $\cF$ we use here is $\cF = \{ g\in \cH_r \text{ s.t. } \|g\|_{\cH_r} \leq 1\}$.

We will focus on the case where $L$ is the squared loss, and we will make the following two assumptions:

\begin{thmappasmp}\label{asmp:yinrkhs}
There exist $f^Y_0,f^Y_1 \in \cH_x$ such that $m_t(x) = \left<f^Y_t,k_x(x,\cdot)\right>_{\cH_x}$, i.e.
the mean potential outcome functions $m_0,m_1$ are in $\cH_x$.   Further assume that $\|f^Y_t\|_{\cH_x} \leq K$.
\end{thmappasmp}

\begin{thmappdef}\label{def:eta}
Define $\eta_{Y_t}(x) :=\sqrt{\int_\cY \left(Y_t - m_t(x)\right)^2 p(Y_t|x)}$. $\eta_{Y_t}(x)$ is the standard deviation of $Y_t|x$.
\end{thmappdef}

\begin{thmappasmp}\label{asmp:etainrkhs}
There exists $f^\eta_0, f^\eta_1 \in \cH_x$ such that $\eta_{Y_t}(x) = \left<f^\eta_t,k_x(x,\cdot)\right>_{\cH_x}$, i.e. the conditional standard deviation functions of $Y_t|x$ are in $\cH_x$. Further assume that $\|f^\eta_t\|_{\cH_x} \leq M$.
\end{thmappasmp}
\begin{thmappasmp}\label{asmp:phiinrkhs}
Let $\Phi: \cX \rightarrow \cY$ be an invertible representation function, and let $\Psi$ be its inverse.
We assume there exists a bounded linear operator $\GP: \cH_r \rightarrow \cH_x$ such that $\left<f^Y_t,k_x(\Psi(r),\cdot)\right>_{\cH_x} = \left<f^Y_t,\GP k_r(r,\cdot)\right>_{\cH_x}$. We further assume that the Hilbert-Schmidt norm (operator norm) $\|\GP\|_{HS}$ of $\GP$ is bounded by $K_\Phi$.
\end{thmappasmp}

The two assumptions above amount to assuming that $\Phi$ can be represented as one-to-one linear map between the two Hilbert spaces $\cH_x$ and $\cH_r$.

Under Assumptions \ref{asmp:yinrkhs} and \ref{asmp:phiinrkhs} about $m_0,m_1$, and $\Phi$, we have that $m_t(\Psi(r)) = \left<\GP^* f^Y_t, k_r(r,\cdot)\right>_{\cH_r}$, where $\GP^*$ is the adjoint operator of $\GP$ \cite{grunewalder2013smooth}.


\begin{thmapplem}\label{lem:mmdbound}
Let $h : \cR \times \{0,1\} \rightarrow \R$ be an hypothesis, and assume that there exist $f^h_t \in \cH_r$ such that $h(r,t) = \left<f^h_t,k_r(r,\cdot)\right>_{\cH_r}$, and such that $\|f^h_t\|_{\cH_r} \leq b$.
Under Assumption \ref{asmp:yinrkhs} about $m_0,m_1$, we have that $\lythr = \int_\cY \left(Y_t - h(r,t)\right)^2 p(Y_t|r) dY_t$ is in the tensor Hilbert space $\cH_r \otimes \cH_r$. Moreover, the norm of $\lythr$ in $\cH_r \otimes \cH_r$ is upper bounded by $4\left(K_\Phi^2 K^2 + b^2\right)$.
\end{thmapplem}
\begin{proof}
We first decompose $\int_\cY \left(Y_t - h(r,t)\right)^2 p(Y_t|x) dY_t$ into a noise and mean fitting term, using $r=\Phi(x)$:
\begin{align}
&\lythr = \nonumber \\
&\int_\cY \left(Y_t - h(r,t)\right)^2 p(Y_t|r) \, dY_t = \nonumber \\
&\int_\cY \left(Y_t - m_t(x) + m_t(x) - h(\Phi(x),t)\right)^2 p(Y_t|x) \, dY_t =\nonumber \\
&\int_\cY \left(Y_t - m_t(x) \right)^2 p(Y_t|x) \, dY_t + \nonumber \\
& \quad  \left(m_t(x) - h(\Phi(x),t)\right)^2  + \nonumber \\
& \quad 2 \int_\cY \left(Y_t - m_t(x) \right)\left(m_t(x) - h(\Phi(x),t)\right) p(Y_t|x) dY_t   =\label{eq:expctzero2} \\
& \eta^2_{Y_t}(x) +  \left(m_t(x) - h(\Phi(x),t)\right)^2  + 0, \label{eq:mmddecomp}
\end{align}
where equality \eqref{eq:expctzero2} is by Definition \ref{def:eta} of $\eta$, and because $\int_\cY \left(Y_t - m_t(x)\right) p(Y_t|x) \, dY_t = 0$ by definition of $m_t(x)$.

Moving to $\cR$, recall that $r=\Phi(x)$, $x = \Psi(r)$.
By linearity of the Hilbert space, we have that $m_t(\Psi(r)) - h(r,t) = \left<\GP^* f^Y_t, k_r(r,\cdot)\right>_{\cH_r} - \left<f^h_t,k_r(r,\cdot)\right>_{\cH_r} = \left<\GP^* f^Y_t - f^h_t,k_r(r,\cdot)\right>_{\cH_r}$.
By a well known result \citep[Theorem 7.25]{steinwart2008support}, the product $(Y_t(\Psi(r)) - h(r,t) ) \cdot (Y_t(\Psi(r)) - h(r,t) )$ lies in the tensor product space  $\cH_r \otimes \cH_r$, and is equal to $\left< (\GP^* f^Y_t - f^h_t) \otimes (\GP^* f^Y_t - f^h_t) , k_r(r,\cdot) \otimes k_r(r,\cdot) \right>_{\cH_r \otimes \cH_r}$. The norm of this function in $\cH_r \otimes \cH_r$ is $\| \GP^* f^Y_t - f^h_t \|^2_{\cH_r}$. This is the general Hilbert space version of the fact that for a vector $w \in \R^d$ one has that $\|w w^\top\|_F = \|w\|_2^2$, where $\|\cdot\|_F$ is the matrix Frobenius norm, and $\| \cdot\|^2_2$ is the square of the standard Euclidean norm.
We therefore have a similar result for $\eta^2_{Y_t}$, using Assumption \ref{asmp:etainrkhs}: $\eta^2_{Y_t}(x) = \eta^2_{Y_t}(\Psi(r)) = \left<\GP^* f^\eta_t \otimes \GP^* f^\eta_t, k_r(r,\cdot) \otimes k_r(r,\cdot) \right>_{\cH_r \otimes \cH_r}$.  
The norm of this function in $\cH_r \otimes \cH_r$ is $\| \GP^* f^\eta_t \|^2_{\cH_r}$. 
Overall this leads us to conclude, using Equation \eqref{eq:mmddecomp} that $\lythr \in  \cH_r \otimes \cH_r$.
Now we have, using \eqref{eq:mmddecomp}:
\begin{align}
&\|\lythr\|_{\cH_r \otimes \cH_r} = \nonumber \\
&\|(\GP^* f^Y_t - f^h_t) \otimes (\GP^* f^Y_t - f^h_t)  + \GP^* f^\eta_t \otimes \GP^* f^\eta_t\|_{\cH_r \otimes \cH_r} \leq \label{eq:triineqmmd} \\
&\| \GP^* f^Y_t - f^h_t \|^2_{\cH_r} + \| \GP^* f^\eta_t \|^2_{\cH_r}  \label{eq:mmdineq1} \leq \\
&2 \|\GP^* f^Y_t\|^2_{\cH_r} + 2\|f^h_t \|^2_{\cH_r}  + \| \GP^* f^\eta_t \|^2_{\cH_r} \label{eq:mmdineq2} \leq \\
&  \|\GP^*\|^2_{HS} \left(2\|f^Y_t\|^2_{\cH_x} + \|f^\eta_t\|^2 _{\cH_x} \right)+ 2\|f^h_t \|^2_{\cH_r} =  \label{eq:mmdeq1}\\
&  \|\GP\|^2_{HS} \left(2\|f^Y_t\|^2_{\cH_x} + \|f^\eta_t\|^2 _{\cH_x} \right)+ 2\|f^h_t \|^2_{\cH_r} \leq \label{eq:mmdineq3}\\
&2 K_\Phi^2 (K^2 + M^2) + 2 b^2. \nonumber
\end{align}
Inequality \eqref{eq:triineqmmd} is by the norms given above and the triangle inequality.
Inequality \eqref{eq:mmdineq1} is because for any Hilbert space $\cH$, $\|a-b\|^2_\cH \leq 2 \|a\|_\cH^2 + 2 \|b\|_\cH^2$. Inequality \eqref{eq:mmdineq2} is by the definition of the operator norm. Equality \eqref{eq:mmdeq1} is because the norm of the adjoint operator is equal to the norm of the original operator, where we abused the notation $\|\cdot \|_{HS}$ to mean both the norm of operators from $\cH_x$ to $\cH_r$ and vice-versa. Finally, inequality \eqref{eq:mmdineq3} is by Assumptions \ref{asmp:yinrkhs}, \ref{asmp:etainrkhs} and \ref{asmp:phiinrkhs}, and by the Lemma's premise on the norm of $f^h_T$.
\end{proof}

\begin{thmapplem}\label{thm:mmdA}
Let $u = p(t=1)$ be the marginal probability of treatment, and assume $0<u<1$. Assume the distribution of $Y_t$ conditioned on $x$ follows Assumptions \ref{asmp:etainrkhs} with constant $M$.
Let $\Phi : \cX \rightarrow \cR$ be a one-to-one representation function which obeys Assumption \ref{asmp:phiinrkhs} with corresponding operator $\GP$ with operator norm $K_\Phi$. Let the functions $Y_0$, $Y_1$ obey Assumption \ref{asmp:yinrkhs}, with bounded Hilbert space norm $K$ . Let $h : \cR \times \{0,1\} \rightarrow \R$ be an hypothesis, and assume that there exist $f^h_t \in \cH_r$ such that $h(r,t) = \left<f^h_t,k_r(r,\cdot)\right>_{\cH_r}$, such that $\|f^h_t\|_{\cH_r} \leq b$. Assume that $\epsilon_F$ and $\epsilon_{CF}$ are defined with respect to $L$ being the squared loss. Then:
\begin{align}\label{eq:thm_mmdA}
&\epsilon_{CF}(h,\Phi) \leq \nonumber\\
&\quad (1-u)  \epsilon^{t=1}_F(h,\Phi) + u  \epsilon^{t=0}_F(h,\Phi)  + \nonumber\\
& \quad 2 \left(K_\Phi^2 (K^2 +M^2) + b^2\right)  \cdot \text{MMD}(\pt_\Phi\, ,  \pc_\Phi),\nonumber\\
\end{align}
where $\epsilon_{CF}$ and $\epsilon_F$ use the squared loss.
\end{thmapplem}
\begin{proof}
We will apply Lemma \ref{lem:genA} with $\cF = {f \in \cH_r \otimes \cH_r \text{ s.t. } \|f\|_{\cH_r \otimes \cH_r} \leq 1}$. By Lemma \ref{lem:mmdbound}, we have that for $B_\Phi = 2\left(K_\Phi^2 (K^2+M^2) + b^2\right)$ and $L$ being the squared loss, $\frac{1}{B_\Phi} \lythr \in \cF$. Inequality \eqref{eq:thm_mmdA} then holds as a special case of Lemma \ref{lem:genA}.
\end{proof}

%\newpage
\begin{thmappthm}
Under the assumptions of Lemma \ref{thm:mmdA}, using the squared loss for $\epsilon_F$, we have:
\begin{align*}
&\epehe(h,\Phi) \leq \nonumber \\
&2 \epsilon^{t=0}_F(h,\Phi)  + 2\epsilon^{t=1}_F(h,\Phi)  - 4\sigma^2_Y +  \nonumber \\
&4 \left(K_\Phi^2 (K^2 +M^2) + b^2\right)  \cdot \text{MMD}( \pt_\Phi \, , \pc_\Phi) .
\end{align*}
\end{thmappthm}
\begin{proof}
Plug in the upper bound of Lemma \ref{thm:mmdA} into the upper bound of Theorem \ref{thm:indtausqloss}.
\end{proof}
